\section{非2的幂次方向量加法实现}

在之前所有的应用程序中，使用的每个向量的元素数量都是 32 的倍数。例如在第一个应用程序中，为 A、B 和 C 这三个向量每个都使用了 1024 个元素。这个值就是 32 的倍数。如果用它除以 32，结果也是 32。在第二个应用程序或者说第二个例子中，每个向量使用了 2048 个元素，这也是 32 的倍数。这样做的原因是任何 CUDA 应用程序的最小规模都是一个 Warp。

\subsection{Warp与线程数的约束}

任何一个 Warp 都由 32 个线程组成。这意味着可以用于任何 CUDA 应用程序的最小线程数是 32 个线程，不能使用例如 16 个线程，必须使用 32 个线程或 32 的倍数。因为必须选择 Warp 的数量，而每个 Warp 由 32 个线程组成。线程总数必须是 32 的倍数。

回答一个可能想到的问题。如果每个向量的向量大小是 16 个元素，这是一个很小的规模。完全可以在 CPU 端相加这些小的规模，而不是在 GPU 端。因为把数据从 CPU 端复制移动到 GPU 端等操作会有开销。

但假设一下，假设向量大小只等于 16 个元素。需要将向量 A 中的这 16 个元素与向量 B 中对应的元素相加，并将结果存储在向量 C 的对应位置。在这种情况下，可以为这个应用程序使用 32 个线程的大小，因为这是最小值。但可以将 16 个元素分配给 16 个线程。而另外 16 个线程不执行任何操作，它们只是处于空闲状态，但在这种情况下，利用率或线程利用率是 50\%。

有 32 个线程，但只使用了其中的 16 个。而这就是不选择 32 的倍数作为元素数量会带来的性能影响之一。现在明白为什么选择每个向量 1024 个元素，而不是例如 1000 个了。这只是本章的引言。必须使用 32 个线程，每个应用程序不能使用少于 32 个线程。因为最小值是一个 Warp，而一个 Warp 等于 32 个线程，每个应用程序必须是 32 或 32 的倍数。这更容易确定 Warp 数量或块大小，因为块大小等于 Warp 数量乘以 32，而每个 Warp 由 32 个线程组成。

当选择块大小时，可以计算网格大小。因为网格大小乘以块大小（以线程计）必须等于元素数量，这不必严格相等，但应该等于向量中元素数量的倍数。因为在后续章节中，将为每个线程分配多个操作。例如这里如果每个向量有 1024 个元素，那么必须选择网格大小和块大小，以提供总共 1024 个线程。在那种情况下，每个元素将被分配给一个不同的线程。

但可以调整应用程序、代码，为每个线程分配两个元素。在那种情况下，每个线程可以相加两个元素，而不仅仅是一个元素。只需要调整代码中的一两行就能做到。并且这将导致减少应用程序中的总块数，而这可能会提高应用程序的占用率（Occupancy）。因为如果每个线程将要执行两个元素，而不是每个元素一个线程，那么将需要总共 512 个线程。如果总共有 512 个线程，而不是 1024 个，那么就可以减少总块数。如果减少了总块数，这意味着在某些情况下可能对性能产生正面或负面的影响。

\subsection{块大小与线程利用率}

那么现在假设有一个不同的例子。假设例如每个向量的元素数量等于 1000，而不是 1024。这里假设了两种情况，如所示，第一种情况是选择 10 个块，其中每个块包含 100 个元素。不是说 100 个线程，因为这是不可能的。每个块将由 96 个线程（等于三个 Warp）或 128 个线程（等于四个 Warp）组成。但在这里，在第一种情况下，只是说每个块将处理 100 个元素。不是说 100 个线程，有 10 个块。其中，每个块将处理 100 个项目，100 个元素。在第二种情况下，选择 8 个块，其中每个块将处理 128 个元素。

这是第一种情况，这是第二种情况。如所示，如果选择每个块 100 个元素，那么需要超过 100 个线程。这里能选择三个 Warp 吗？三个 Warp 将等于 96 个线程，这行不通。必须选择下一个 32 的倍数，下一个 32 的倍数是 128。

正如在这里看到的，例如在块内，必须选择四个 Warp。所有这些 Warp 的线程总数将是 128 个线程。只需要大约 100 个，这是 32、64、96。将只使用最后一个 Warp 中的 4 个线程，这里使用的线程总数是每个块 100 个线程。而有多少？有 28 个线程处于空闲状态。这适用于所有 10 个块。

如所示，每个块使用了完整的三个 Warp 以及最后一个 Warp 的 4 个线程。这里可以计算线程利用率，线程利用率将是 78\%。因为如果用使用的线程数 100 除以 128，这将给出 78\%。即 28 个空闲线程除以线程总数，这是使用的部分，这是从所有可用线程中利用的部分。

在第二种情况下，减少了块数。为什么减少了块数？因为每个块将处理 128 个元素，而只有 1000 个，所以不需要使用 10 个块。8 个块就足够了。将使用所有四个 Warp，如在此所示。

总之在第二种情况下，如在这里看到的，例如对于块内，将使用 128 个线程。同样的情况适用于第二个块，也适用于第三个块。以此类推，直到最后一个块。

在最后一个块中将使用四个 Warp 中的三个和最后一个 Warp 的部分线程。只有这些。能告诉所有这些块中有多少个线程处于空闲状态吗？它们大约是 32 减 8，这大约是 24 个线程处于空闲状态。但在这里每个块大约有 28 个线程处于空闲状态（指第一种情况），不是所有块的总和。几乎是 100\% 利用率（指第二种情况的前 7 个块）。这就是为什么这里的平均线程利用率是 96.6\%。

现在看到选择块数量和每个块的线程数的影响了吗？这非常重要。在这个部分（第二种情况）使用了更多线程，因为这在线程资源利用率方面会更好。而这里（第一种情况）有更多的线程处于空闲状态，同时还减少了块的数量（第二种情况），这也不错。因为这里每个块的大小和那里每个块的大小是恒定的。这里使用了四个 Warp，那里也使用了四个 Warp。但区别是，这里利用了所有线程或所有 Warp，而那里只利用了其中一部分，而且有 10 个块，而这里有 8 个块。

\subsection{波次与执行效率}

假设 GPU——这只是个假设，因为这是个很小的数字——假设 GPU 包含大约 8 个 SM。只有 8 个 SM。

在第二种情况下，将每个 SM 分配一个块，并且可以在一个波次（Wave）内完成执行。这里的波次执行前 8 个块。但在那里（第一种情况）必须进行两个波次，而不是一个。因为需要执行 10 个块，每个块将被分配给一个 SM，因为只有 8 个 SM。有两个波次，这将使执行时间加倍。然后在第二个波次中，其中 2 个 SM 将执行剩余的 2 个块。

这只是一个例子，展示当每个向量的元素数量不是 32 的倍数时该怎么办——每个块的线程数，每个网格的块数等。希望能自己应用这个。因为已经多次解释了，唯一需要做的就是控制元素数量。这就是需要做的全部：相同的代码，相同的核函数，一切都相同。唯一需要操纵的就是这三个数字。
